\subsection{The specified lifts do not have specified tripod values}

\statement{B.42}{Proved}{A nontrivial arithmetic ambiguity in every E.58 lift}
Put $H_0=F'F^3\trianglelefteq\widehat B$ and
\[
 R_{\rm sol}=\bigcap_{G_{\mathbb Q}\twoheadrightarrow D,
                         \ D\ \mathrm{finite\ solvable}}\ker(G_{\mathbb Q}\to D).
\]
This is a nontrivial closed subgroup of $G_{\mathbb Q}$.
For every $\sigma\in R_{\rm sol}$, the normalized modular
automorphism induced by $\sigma$ has multiplier one and is
the identity on $\widehat B/R_{\rm nil}(H_0)$.  Its two
canonical pure accessory frames have image one in this
same relative quotient.  Consequently right multiplication
by any such $\sigma$ preserves every prescribed relative
automorphism and relative frame coordinate in the
construction of LM7\_PUSH E.58, including its finite
modulo-twelve frame classes.

\begin{proof}
The cover of the thrice-punctured line associated with
$H_0\cap F=H_0$ is
\[
 E^\circ:\quad X^3+Y^3=Z^3,\quad XYZ\ne0,
       \qquad (X:Y:Z)\longmapsto X^3/Z^3.
                                                               \tag{B.42.1}
\]
Its degree is nine, its compactification has genus one,
and it has nine punctures.  These assertions follow
either from this smooth plane cubic or from the three
ramification indices of the $(\mathbb Z/3)^2$-cover.
All nine cusps and the automorphism
$(X:Y:Z)\mapsto(\zeta_3X:Y:Z)$ are defined over
$k=\mathbb Q(\zeta_3)$.  The rational cusp $(0:1:1)$
is fixed by this automorphism and can be used as origin.
Its action on the compact Tate module has minimal
polynomial $T^2+T+1$, so its centralizer in
$\operatorname{GL}_2(\mathbb Q_p)$ is a commutative
quadratic algebra, including at $p=2,3$.

The standard tangential Galois action on $F$ restricts
on this cover to its arithmetic action.  One can check
the comparison directly near $0$: choose the Puiseux
branches $X/Z=t^{1/3}$ and
$Y/Z=(1-t)^{1/3}$ with leading coefficient one.
After the indicated cyclotomic extension all branches
and cusp labels are fixed.  Changing the lifted
tangential point changes the action by a geometric
inner automorphism and does not change homology.
Thus on $G_k$ the exact homology sequence is
\[
 0\longrightarrow\mathbb Z_p(1)^8
 \longrightarrow H_1(E^\circ_{\overline k},\mathbb Z_p)
 \longrightarrow T_pE\longrightarrow0.                     \tag{B.42.2}
\]
The first diagonal action is scalar cyclotomic and
the last lies in the commutative centralizer just
described.  The kernel of the action on these two
diagonal blocks is a subgroup of the additive
module $\operatorname{Hom}_{\mathbb Z_p}
(T_pE,\mathbb Z_p(1)^8)$.  Therefore the arithmetic
image on (B.42.2) is prosolvable.  Passing from $G_k$
to $G_{\mathbb Q}$ adds a quotient of order at most
two and preserves that conclusion.  No assertion
that the genus-twenty-five level-twelve Jacobian
is CM is needed or made.

Let $P_p=\widehat B/O^p(H_0)$ and
$U_p=H_0^{(p)}\trianglelefteq P_p$.
Here $U_p$ is free pro-$p$ of rank ten and
$P_p/U_p=\widehat B/H_0$ has order fifty-four.
Conjugation by this finite deck group on
$H_1(U_p,\mathbb Z_p)$ is faithful.  Indeed its
boundary module records its permutation of the
nine punctures faithfully, and an automorphism
of a complex elliptic curve fixing nine distinct
points is the identity.  To verify the latter
elementary bound, write such an automorphism as
$z\mapsto az+b$ on a complex torus.  A nontrivial
translation has no fixed point; when $a\ne1$,
the possible elliptic units give at most four
fixed points.  It follows that
$C_{P_p}(U_p)=1$: this centralizer meets $U_p$
in its trivial center and its deck image acts
trivially on the displayed homology.

Restriction consequently embeds
$\operatorname{Aut}(P_p;U_p)$ in
$\operatorname{Aut}(U_p)$.  An automorphism
restricting to the identity differs from the
identity on every coset representative by an
element of $C_{P_p}(U_p)$, which proves this
injectivity.  The kernel of
\[
 \operatorname{Aut}(U_p)\longrightarrow
     \operatorname{GL}(U_p/\Phi(U_p))
\]
is pro-$p$.  For clarity this is the full
automorphism group of a finitely generated
free pro-$p$ group: on its finite characteristic
$p$-group quotients the automorphisms trivial
on the Frattini quotient form $p$-groups,
as successive Frattini-series corrections
embed in additive $\mathbb F_p$-modules;
the inverse limit proves the assertion.
Combined with (B.42.2), this proves that the
actual arithmetic image in
$\operatorname{Aut}(P_p;U_p)$ is prosolvable.
Also $P_p$ itself is prosolvable, being pro-$p$
by the solvable group of order fifty-four.

We must retain the frames as well as the
automorphisms: identity of the latter alone
does not imply identity of relative frames,
as E.58 itself illustrates.  On $G_k$, the
canonical normalized frames $\beta_g,A_g\in F$
are unique and satisfy
\[
 \beta_{gh}=\Phi_g(\beta_h)\beta_g,\qquad
        A_{gh}=\Phi_g(A_h)A_g.                              \tag{B.42.3}
\]
Here $g$ has multiplier congruent to one
modulo three, so the order-three generator
is sent to a conjugate of itself.  The
formula follows by composing the two actual
automorphisms and using the trivial pure
centralizers of $s$ and $r$ in the free
profinite product.  Thus
\[
 g\longmapsto
  (\beta_g^{-1},A_g^{-1},\Phi_g)
\]
is a continuous homomorphism into
$(P_p\times P_p)\rtimes\rho_p(G_k)$ after
projection to $P_p$.  That group is
prosolvable by the preceding paragraphs.

The prosolvable residual of $G_k$ equals
$R_{\rm sol}$.  To see the nontrivial direction,
take a finite solvable quotient of the
normal index-two subgroup $G_k$ and
intersect its kernel with its conjugate
under $G_{\mathbb Q}$.  The resulting
quotient of $G_{\mathbb Q}$ is an extension
of a subgroup of a product of solvable
groups by $C_2$, and is solvable.  It
follows that every $\sigma\in R_{\rm sol}$
kills the framed representation above,
for every $p$.  The actual fibre-product
identity for $\widehat B/R_{\rm nil}(H_0)$
then gives the assertion for that whole
relative pronilpotent quotient.  The
cyclotomic character is abelian and also
vanishes on $R_{\rm sol}$.

The modulo-twelve frames really are retained by
these relative data.  With the global projective
sign fixed by its modulo-four value, the exact
Chinese remainder description is
\[
 q_{12}(F)=V_4\times\operatorname{SL}_2(\mathbb F_3),
       \qquad q_{12}(H_0)=V_4\times Q_8.
\]
Here $V_4$ is the complement to $\langle-I\rangle$
in the modulo-four kernel, represented by
$I$, $\left(\begin{smallmatrix}1&2\\0&1\end{smallmatrix}\right)$,
$\left(\begin{smallmatrix}1&0\\2&1\end{smallmatrix}\right)$
and $\left(\begin{smallmatrix}1&2\\2&1\end{smallmatrix}\right)$.
The second equality follows because the derived
subgroup of $\operatorname{SL}_2(\mathbb F_3)$
is $Q_8$, its quotient has order three, and
cubing preserves $V_4$.  This is the same
marked finite quotient as in the refereed R2,
not an unsigned replacement of it.  Since
$q_{12}(H_0)$ is a $2$-group, $q_{12}$ factors
through $\widehat B/O^2(H_0)$, as required.

Finally $R_{\rm sol}\ne1$.  The polynomial
$x^5-x-1$ is irreducible modulo three,
and its square-free factorization modulo
two has degrees two and three.  Its
Galois group is therefore transitive of
prime degree and contains a transposition
(the cube of that Frobenius permutation).
A primitive permutation group containing
a transposition is the full symmetric
group: the graph of conjugates of that
transposition is invariant and connected,
and its edge transpositions generate the
symmetric group.  Thus this splitting
field has group $S_5$.  The image of
$R_{\rm sol}$ in it is $A_5$, since the
remaining quotient is solvable and the
sign quotient already is solvable.  For
example there exists $\sigma\in R_{\rm sol}$
whose image is a prescribed three-cycle.
The two finite-field checks are included
in \texttt{check\_lift\_ambiguity.py}.
\end{proof}

\statement{B.43}{Proved}{A finite simple group excluded uniformly from these relative levels}
Let $S=A_{10000}$.  No subgroup of
$\operatorname{GL}_d(\mathbb F_p)$, for $d\leq97$
and any prime $p$, has $S$ as a quotient of a
subgroup.  Consequently every finite section
of a closed subgroup of
$\operatorname{Aut}(\widehat F_d^{(p)})$
also avoids $S$.

\begin{proof}
Pad a representation with trivial summands
to dimension $97$.  For a prime $q\ne p$,
a finite $q$-subgroup $Q$ of
$\operatorname{GL}_{97}(\mathbb F_p)$ has
a faithful characteristic-zero representation
of the same dimension.  One direct verification
is to lift its matrices successively modulo
$p^r$.  The obstruction at the next step is
a $2$-cocycle with coefficients in the
matrix module over $\mathbb F_p$; its
cohomology vanishes by averaging over $Q$,
whose order is prime to $p$.  Corrections
therefore give compatible homomorphisms
to $\operatorname{GL}_{97}(\mathbb Z/p^r)$.
Their limit is faithful because its reduction
is faithful.  A finite-group representation
in characteristic zero has the same character
as a complex representation, giving the claim.

Collins's Theorem~A, \cite{CECollinsJordan2007},
now gives an abelian normal subgroup
$A\trianglelefteq Q$ of index at most $98!$.
This is an application in characteristic
zero to $Q$, not an assertion of Jordan's
theorem for arbitrary groups in
characteristic $p$.  The abelian group
$A$ is simultaneously diagonalizable
in dimension $97$, so every subgroup
and every quotient of it needs at most
$97$ generators.

If an elementary abelian group $C_q^r$
is a section $L/M$ of $Q$, the image
of $A\cap L$ in $L/M$ has rank at most
$97$.  The remaining quotient has order
at most $[Q:A]\leq98!$.  Therefore
\[
 r\leq97+\lfloor\log_q(98!)\rfloor
   =\begin{cases}608,&q=2,\\419,&q=3.\end{cases}            \tag{B.43.1}
\]
Now $A_{10000}$ contains $C_2^{2500}$,
using a double transposition on each
of $2500$ disjoint four-letter blocks,
and contains $C_3^{3333}$, using
disjoint three-cycles.  If
$L/M\simeq A_{10000}$ with
$L\leq\operatorname{GL}_{97}(\mathbb F_p)$,
take the inverse image in $L$ of one
of these elementary abelian subgroups.
A $q$-Sylow subgroup of that inverse
image maps onto it.  Choose $q=2$
when $p\ne2$, and $q=3$ when $p=2$.
In each case (B.43.1) is contradicted.
This proves the assertion about all
sections, not only faithful embeddings
of the simple group.

The reduction of
$\operatorname{Aut}(\widehat F_d^{(p)})$
on the Frattini quotient has pro-$p$
kernel, as proved in B.42.  The class
of finite groups with no $S$-section
is closed under subgroups, quotients
and extensions.  For the extension
assertion, the image in a putative
simple quotient $S$ of the intersection
with a normal subgroup is either $1$
or $S$; the two cases put an $S$-section
in the quotient or in the normal
subgroup, respectively.  Applying
this observation to the Frattini
reduction proves the last assertion.
\end{proof}

\statement{B.44}{Proved}{A simultaneous nontrivial arithmetic kernel for both lifting data}
Let $N_{12}=\ker(\widehat B\to
\operatorname{SL}_2(\mathbb Z/12)/\{\pm I\})$.
There is a nontrivial closed subgroup
$Z\leq G_{\mathbb Q}$ with multiplier
one whose normalized automorphisms
and both canonical accessory frames
are the identity at every relative
pro-$p$ level over $H_0$ and over
$N_{12}$, simultaneously.  Thus it
kills the whole relative pronilpotent
source and the specified relative
frame coordinates in both LM7\_PUSH
E.58 and LM7 D.46.

\begin{proof}
We keep the source and frame data,
not just the images in graded Lie
algebras.  The two open torsion-free
subgroups have indices $54$ and
$576$ in $\widehat B$, and free
ranks $10$ and $97$.  The first
index is $6\cdot9$; the second is
$|\operatorname{SL}_2(\mathbb Z/4)|
 |\operatorname{SL}_2(\mathbb F_3)|/2
=48\cdot24/2$.  Both finite
quotients embed the two torsion
factors, so their kernels are
torsion-free.  Multiplicativity
of the Euler characteristic
$\chi(C_2*C_3)=-1/6$ gives the
stated ranks.  The normalized
Galois action preserves $H_0$
because it is characteristic.
It preserves $N_{12}$ because
the level-twelve torsion cover
of the modular curve is arithmetic:
the Galois action on level structures
lies in $\operatorname{GL}_2(\mathbb Z/12)$
and normalizes the geometric
congruence kernel.  Equivalently
this is the actual congruence
covariance in the refereed R3.

For either subgroup $U$ put
$P_{U,p}=\widehat B/O^p(U)$ and
$U_p=U^{(p)}$.  Thus
\[
 1\longrightarrow U_p\longrightarrow
 P_{U,p}\longrightarrow D_U
 \longrightarrow1,
 \quad |D_U|\leq576,
 \quad \operatorname{rank}(U_p)\leq97.                       \tag{B.44.1}
\]
Let $\mathscr C$ be the class of
finite groups with no $A_{10000}$
section.  Besides its elementary
closure properties proved in B.43,
the class of pro-$\mathscr C$
groups is closed under arbitrary
products and closed subgroups.
Indeed a continuous finite quotient
of a closed subgroup of a product
factors through a subgroup of a
finite product of finite quotients:
choose a product-open normal
neighborhood whose intersection
lies in the given finite-index
kernel.  The finite product case
follows from extension closure.
The same argument, or passage
to finite quotients in an exact
sequence, proves extension closure
for pro-$\mathscr C$ groups.

The group $P_{U,p}$ is pro-$\mathscr C$
by (B.44.1).  So is
$\operatorname{Aut}(P_{U,p};U_p)$.
To check the latter without a
faithful-deck-action assumption,
restriction to $U_p$ has finite
kernel of order at most $576^{576}$.
Indeed $C_{P_{U,p}}(U_p)$ meets
$U_p$ trivially, since a free
pro-$p$ group of rank at least
two has trivial center, and hence
has at most $|D_U|$ elements.
An automorphism fixing $U_p$
pointwise is determined by its
values on $|D_U|$ coset representatives;
each value differs from the
original representative by an
element of that centralizer.
This proves the bound.  The
integer $576^{576}$ is smaller
than $|A_{10000}|=10000!/2$,
so that finite kernel is in
$\mathscr C$.  The restriction
image is pro-$\mathscr C$ by
B.43.  Extension closure proves
the assertion.  This bound also
covers every prime $2,3,5$.

On $G_{\mathbb Q(\zeta_3)}$, use
the actual framed homomorphisms
of (B.42.3) to
\[
 (P_{U,p}\times P_{U,p})
   \rtimes\operatorname{Aut}(P_{U,p};U_p),
             \qquad U\in\{H_0,N_{12}\}.
\]
Take their product over both $U$
and all primes, and also retain
the full cyclotomic character.
The resulting continuous
homomorphism $\Theta$ has a
pro-$\mathscr C$ image.  Put
$Z=\ker\Theta$.  By construction
it kills all the asserted
automorphisms and relative frames
and has multiplier one.  It
remains to prove that it is
nontrivial, rather than infer
that from the sizes of some
finite images.

Here is a completely specified
large symmetric Galois quotient.
Let $n=10000$ and
\[
                    f_n(x)=x^n+15x^{n-1}+10x^2+30.          \tag{B.44.2}
\]
It is Eisenstein at five, so
its Galois group is transitive.
Modulo two it factors as
$x^{n-1}(x+1)$.  Hensel lifting
the two coprime factors gives
a factor of degree $n-1$ whose
constant term has valuation one
and whose other nonleading
coefficients are divisible by
two.  That factor is Eisenstein;
since $n-1=9999$ is odd, its
totally ramified extension is
tame.  Tame inertia in the
Galois closure is cyclic and
acts transitively on these
$n-1$ roots; the remaining
linear root is unramified.
Thus the Galois group contains
an $(n-1)$-cycle.

Modulo three the factorization
is $x^2(x^{n-2}+1)$.  The latter
factor is separable and is a
unit at zero, since $n-2=9998$
is nonzero modulo three.
Hensel lifting gives an
Eisenstein quadratic factor
and an unramified splitting
field for all remaining roots.
Its tame quadratic inertia
therefore gives one transposition.
A transitive group containing
an $(n-1)$-cycle is doubly
transitive: the stabilizer of
the fixed letter of that cycle
is transitive on the other
letters.  It is therefore
primitive, and the transposition
argument in B.42 gives $S_n$.
These elementary local conditions
are also the specialization
$p=2,q=3,r=5$ of Osada's
Corollary~1 to Theorem~3,
\cite{CEOsada1987II}; we have
included the local argument
rather than invoking an
unspecified inverse-Galois
realization.

On $G_{\mathbb Q(\zeta_3)}$
the image of this Galois
quotient is $S_n$ or $A_n$,
since the base change has
degree two.  If $Z$ had
trivial image there, that
quotient would factor through
$\operatorname{im}\Theta$,
contrary to its pro-$\mathscr C$
property.  Its nontrivial
normal image therefore contains
$A_n$.  In particular there
exists $\sigma\in Z$ mapping
to a prescribed three-cycle.
This proves $Z\ne1$ through
one actual finite Galois
quotient.  The proof does not
compute its splitting field
numerically.  The exact
factorization conditions and
all integer inequalities used
above are checked by
\texttt{check\_lift\_ambiguity.py}.
\end{proof}

\statement{B.45}{Proved}{The relative prescriptions leave all tripod and corrected resolvent values undetermined}
Fix any actual full lift $\Phi$
allowed by LM7\_PUSH E.58, or
by R2 applied to the specified
relative action and frames of
LM7 D.46.  For the subgroup
$Z$ of B.44, every $\Phi\sigma$,
$\sigma\in Z$, is another
allowed lift with the same
relative data.  For every
three-subset $t$ of the
icosahedral cusps, the map
\[
 Z\longrightarrow\mathcal P,
       \qquad \sigma\longmapsto
                         \mathfrak T_t(\Phi\sigma)           \tag{B.45.1}
\]
is injective.  The same is true
of $\sigma\mapsto\mathfrak R(\Phi\sigma)$
for each of the three typed
resolvent homomorphisms
$\mathfrak R$ constructed in B.53:
the first two have target
$\widehat{GT}_0$, and the third
has the peripheral outer
automorphism group of
$C_3\amalg C_3$ as its target.
Thus neither existential
construction specifies a single
value of any tripod test or
of any of these corrected
resolvent maps.
At most one element of the
coset $\Phi Z$ can belong
to $\mathcal K$.

\begin{proof}
Each lift preserves its
specified subgroup $U$ and
all of its characteristic
pro-$p$ residuals, because
it induces the given
automorphism of the relative
quotient over $\widehat B/U$.
The action of $\sigma$ and
both its frames have image
one in every $P_{U,p}$.
Composition (B.42.3) therefore
leaves the relative frames
of $\Phi$ unchanged, as
well as its automorphism.
Its multiplier remains one,
so $\Phi\sigma\in\mathcal H$.
For E.58, one may use the
larger subgroup $R_{\rm sol}$
of B.42 instead of $Z$.

The full multiplier-one
condition implies that
$\sigma$ fixes $\mathbb Q(\zeta_5)$.
By the properly typed
arithmetic comparison in
B.10, $\mathfrak T_t(\sigma)$
is its usual faithful
Galois action on the
three-punctured sphere.
Consequently
\[
 \mathfrak T_t(\Phi\sigma)
       =\mathfrak T_t(\Phi)\mathfrak T_t(\sigma)
\]
has distinct values for
distinct $\sigma$, proving
(B.45.1).

The lifts $\Phi\sigma$ all have
multiplier one, and hence lie
in the common domain
$\ker\chi_{-15}$ of B.53.
Put $k=\mathbb Q(\sqrt{-15})$.
The full cyclotomic character
of every $\sigma\in Z$ is one,
so $Z\leq G_{k(\zeta_3)}$.
The arithmetic identity in
B.53 gives
$\mathfrak R(\sigma)=\sigma$
for either of its first two
maps.  Since they are actual
homomorphisms, multiplication
by $\mathfrak R(\Phi)$ proves
the asserted injectivity for
these two values as well.

For completeness the third
target also detects $Z$
faithfully, although it is
not $\widehat{GT}_0$.
Use the geometric marking of
B.53 to write the surviving
root stack as the sphere
with cone orders three at
$0,1$ and a cusp at infinity.
Its geometric group is
$D=C_3\amalg C_3$, with positive
cone generators $a,b$.  The
degree-three cover
\[
 \mathbb P^1_z\setminus\mu_3
       \longrightarrow \mathbb P^1_t(3,3;\infty),
 \qquad t=\frac{z^3}{z^3-1},                               \tag{B.45.2}
\]
is finite etale as a map to
the root stack.  Its geometric
fundamental group $P\cong\widehat F_2$
is the kernel of the continuous
epimorphism $D\to C_3$ sending
$a$ to $1$ and $b$ to $-1$.
The cover and the point $z=2$
over $t=8/7$ are defined over
$k$, so their based Galois
actions give compatible genuine
actions on $P\leq D$.

Suppose that $\sigma\in G_{k(\zeta_3)}$
has trivial outer action on $D$.
Its based action is then
conjugation by some $d\in D$.
On $P$ the resulting outer
class is the deck class of
$dP\in D/P=C_3$.  A nontrivial
such deck class cycles the
three cusps $\mu_3$ of the cover.
The Galois action of $\sigma$
fixes all three cusp classes,
so this deck class is trivial.
Consequently $\sigma$ acts
trivially in $\operatorname{Out}(P)$.
After a fractional linear
change of coordinate over
$k(\zeta_3)$, this is the
faithful Galois action on
$\mathbb P^1\setminus\{0,1,\infty\}$
used in B.10.  Thus $\sigma=1$.
In particular the third
resolvent is injective on $Z$;
its homomorphism law again
gives injectivity on the
coset $\Phi Z$.

If two elements
of $\Phi Z$ belonged to
$\mathcal K$, their quotient
would lie in
$\mathcal K\cap G_{\mathbb Q}=1$,
again by B.10.
\end{proof}

This result does not assert
that one member of either
coset actually lies in
$\mathcal K$.  Uniform
exclusion of all admissible
lifts could still be true;
it is not established by
the nonuniqueness argument.
Nor does this identify the
common kernel of the three
resolvents with $\mathcal K$,
or prove the remaining
claims of P5--P7 discussed
in B.12.  It proves that
the specified relative data
cannot define one full
element with determined
tripod or corrected resolvent
values.
